\section{Introdu\c{c}\~ao}\label{sec:intro}

Reformas da jornada legal alteram simultaneamente a utiliza\c{c}\~ao do trabalho, a produtividade por hora e os incentivos de formaliza\c{c}\~ao. Essa intera\c{c}\~ao \'e especialmente relevante no Brasil, onde grande parcela do emprego n\~ao est\'a diretamente coberta pelas restri\c{c}\~oes aplic\'aveis ao setor formal. O debate legislativo de 2025 e abril de 2026 tamb\'em deixou de se resumir a uma escolha bin\'aria entre 44 e 36 horas: passaram a coexistir propostas de 40 horas, escala 5x2 e trajet\'orias faseadas \citep{CamaraPEC82025Protocol,CamaraFortyHours2025,CamaraFortyHoursBill2026}. A quest\~ao quantitativa deste artigo \'e: qual ganho de PTF manteria o produto agregado constante em cada ponto dessa trajet\'oria?

Para respond\^e-la, constru\'imos um \'indice de produtividade requerida, $\Areq$: a varia\c{c}\~ao \'unica de PTF que recomp\~oe exatamente o produto de curto prazo ap\'os a redu\c{c}\~ao do teto. O objeto n\~ao \'e uma previs\~ao de produtividade, mas um contrafactual de neutralidade que permite comparar desenhos de reforma e coloc\'a-los na escala da experi\^encia macroecon\^omica brasileira.

Nosso modelo combina quatro ingredientes. Primeiro, um agregador CES incorpora substitui\c{c}\~ao imperfeita entre trabalho efetivo formal e informal. Segundo, a produtividade hor\'aria varia com a jornada, com pico de efici\^encia assumido em 40 horas. Terceiro, cunhas e custos de ajuste governam a composi\c{c}\~ao formal--informal, com uma extens\~ao setorial. Quarto, prefer\^encias Greenwood--Hercowitz--Huffman (GHH) fornecem um diagn\'ostico transparente de consumo e lazer. Com a elasticidade de substitui\c{c}\~ao fixada em $1{,}326$, a raz\~ao de remunera\c{c}\~ao hor\'aria da PNAD Cont\'inua de 2024T4 disciplina o peso tecnol\'ogico formal. Os resultados tamb\'em cobrem sensibilidades de par\^ametros.

Os resultados revelam uma n\~ao linearidade substantiva. A base limita previamente as horas habituais formais da PNAD a 44h, como aproxima\c{c}\~ao para avaliar as redu\c{c}\~oes do teto, sem trat\'a-las como horas contratadas observadas. A alternativa de 40 horas requer $1{,}53$--$1{,}57\%$ de PTF. No teto de 36 horas, $\Areq=6{,}52\%$ quando a fun\c{c}\~ao de efici\^encia penaliza apenas jornadas acima de 40 horas e $8{,}38\%$ quando a penalidade \'e bilateral. A informalidade aumenta entre $2{,}28$ e $2{,}92$ pontos percentuais, mas a decomposi\c{c}\~ao mostra que a redu\c{c}\~ao f\'isica das horas formais \'e a maior parcela negativa da varia\c{c}\~ao do produto. A magnitude depende da efici\^encia e da distribui\c{c}\~ao inicial: preservar a cauda habitual acima de 44h reduz o requisito para 40 horas a cerca de $0{,}07\%$.

O artigo contribui em tr\^es dimens\~oes. Em primeiro lugar, aproxima a
literatura sobre regula\c{c}\~ao do tempo de trabalho da evid\^encia brasileira sobre jornada, emprego e
produtividade. \citet{Fernandes1991} mostrou que os efeitos de uma redu\c{c}\~ao
generalizada da jornada n\~ao podem ser inferidos diretamente da resposta de
uma firma isolada, enquanto \citet{GonzagaMenezesFilhoCamargo2003} estudam a
redu\c{c}\~ao constitucional de 48 para 44 horas e documentam seus efeitos
sobre jornada efetiva, emprego e sal\'arios. Em uma perspectiva agregada,
\citet{BarbosaFilhoPessoa2014} mostram que medir o insumo trabalho pelo total
de horas, em vez de apenas pelo n\'umero de pessoas ocupadas, altera
substantivamente a leitura da produtividade brasileira. Este artigo
complementa esses trabalhos ao construir um contrafactual prospectivo que
relaciona diferentes tetos legais \`a efici\^encia hor\'aria e ao ganho de PTF
necess\'ario para preservar o produto.

Em segundo lugar, o artigo incorpora explicitamente a margem
formal--informal. A literatura brasileira mostra que formalidade e
informalidade n\~ao devem ser interpretadas como estados isolados ou como uma
simples diferen\c{c}a contratual. Os diferenciais observados de remunera\c{c}\~ao
refletem sele\c{c}\~ao dos trabalhadores
\citep{MenezesFilhoMendesAlmeida2004}, e as transi\c{c}\~oes entre emprego
formal, informalidade, trabalho por conta pr\'opria e desemprego s\~ao
quantitativamente relevantes \citep{CuriMenezesFilho2006,HirataMachado2010}.
Ao mesmo tempo, institui\c{c}\~oes trabalhistas e mecanismos de
fiscaliza\c{c}\~ao afetam conjuntamente formaliza\c{c}\~ao, desemprego e
bem-estar \citep{Ulyssea2008,Ulyssea2018}, enquanto a din\^amica do emprego
formal tamb\'em est\'a associada ao tamanho dos estabelecimentos
\citep{CorseuilMouraRamos2011}. O modelo incorpora essa margem por meio da
substitui\c{c}\~ao imperfeita entre trabalho formal e informal e de cunhas de
formaliza\c{c}\~ao. Isso permite avaliar n\~ao apenas o efeito agregado da reforma, mas tamb\'em como as respostas de produto e composi\c{c}\~ao do emprego diferem entre setores.

Em terceiro lugar, o artigo organiza os contrafactuais em uma curva completa
de $\Areq$. Em vez de avaliar somente um ponto final, o exerc\'icio acompanha
toda a trajet\'oria entre 44 e 36 horas e identifica a n\~ao linearidade
associada \`a efici\^encia hor\'aria. Essa representa\c{c}\~ao permite
distinguir reformas que predominantemente eliminam horas associadas \`a fadiga
daquelas que passam a comprimir horas tratadas como efetivas pelo modelo. A
decomposi\c{c}\~ao entre perda mec\^anica de horas, resposta da efici\^encia e
realoca\c{c}\~ao formal--informal fornece ainda uma interpreta\c{c}\~ao
econ\^omica dos resultados: a redu\c{c}\~ao de horas formais domina as parcelas negativas da resposta do produto agregado; a resposta da efici\^encia depende da especifica\c{c}\~ao, e a realoca\c{c}\~ao para a informalidade amplia a perda.

A op\c{c}\~ao por um benchmark est\'atico de curto prazo, com capital
predeterminado e emprego agregado mantido no n\'ivel inicial, torna transparente
a interpreta\c{c}\~ao de $\Areq$: trata-se da produtividade necess\'aria para
compensar a redu\c{c}\~ao de horas antes que outros mecanismos de ajuste tenham
se completado. As reformas do Brasil em 1988 e de Portugal em 1996 fornecem
contexto externo para interpretar as margens de resposta
\citep{GonzagaMenezesFilhoCamargo2003,AsaiLopesTondini2024}. Em conjunto com
as sensibilidades de par\^ametros, essas compara\c{c}\~oes permitem apresentar os
resultados como refer\^encias quantitativas condicionais para o desenho e o
sequenciamento de reformas da jornada no Brasil.

O restante do artigo est\'a organizado da seguinte forma. A Se\c{c}\~ao~\ref{sec:facts} apresenta os fatos que disciplinam o exerc\'icio. A Se\c{c}\~ao~\ref{sec:model} descreve o modelo. A Se\c{c}\~ao~\ref{sec:calibration} discute calibra\c{c}\~ao e identifica\c{c}\~ao; a Se\c{c}\~ao~\ref{sec:validation} apresenta checagens externas. A Se\c{c}\~ao~\ref{sec:results} reporta os contrafactuais e a Se\c{c}\~ao~\ref{sec:policy} examina implica\c{c}\~oes de transi\c{c}\~ao. A Se\c{c}\~ao~\ref{sec:conclusion} conclui.

% ==================================================================
% 2. Fatos brasileiros
% ==================================================================
